\section{Asymptotic Analysis}
Algorithmic Analysis involves the \textcolor{Bittersweet}{resources needed} by
a algorithm:
\begin{itemize}
  \item {\textbf{Time}}: the number of steps taken to complete. 
  \item {\textbf{Space}}: the amount of memory used to complete. 
    data fits in memory)
\end{itemize}
\textcolor{Blue}{Running time} is a function of the
\textcolor{Bittersweet}{size of the input}, denoted as $n$, which can be
measured in various ways, such as the number of elements in an array, the
number of nodes in a graph, etc.

Typically, we perform \textbf{worst-case analysis} considering the
\textcolor{Bittersweet}{maximum running time} for any input of size (at most)
$n$, or \textbf{average-case analysis} considering the
\textcolor{Bittersweet}{expected running time} over all inputs of size $n$
under some probability distribution.

\paragraph{RAM Model of Computation}
In our model, we analyse algorithms under the assumptions that:
\begin{itemize}
  \item Each basic unit of memory is a \textcolor{Blue}{word}.
  \item \textcolor{Blue}{RAM} is an array of words, storing instructions and
    data.
  \item One unit of time taken to access any word.
  \item \textcolor{Bittersweet}{Constant time} for each operation. (\emph{Except exponentiation})
  \item If input or output is too large, time taken depends on the number of
    words needed to store it.
\end{itemize}
Running time is \textcolor{Bittersweet}{roughly proportional to the number of
instructions} executed in the RAM model, independent of the actual machine.

\subsection{Asymptotic Notation}
\paragraph{Definitions}
\textcolor{Blue}{$\mathbf{f \in O(g)}$} if 
\[\exists c,n_0>0\quad \forall n\ge n_0\quad 0\le f(n)\le c\cdot g(n).\] 
\[\lim_{n \to \infty} \frac{f(n)}{g(n)}<\infty\implies f\in O(g).\] 

$f$ grows \textcolor{Bittersweet}{at most as fast as $g$} asymptotically.

\textcolor{Blue}{$\mathbf{f \in \Omega(g)}$} if
\[\exists c,n_0>0\quad \forall n\ge n_0\quad 0\le c\cdot g(n)\le f(n).\]
\[0<\lim_{n \to \infty} \frac{f(n)}{g(n)}\implies f\in \Omega(g).\]

$f$ grows \textcolor{Bittersweet}{at least as fast as $g$} asymptotically.

\textcolor{Blue}{$\mathbf{f \in \Theta(g)}$} if $f\in O(g)$ and $f\in \Omega(g)$, i.e. 
\[\exists c_1,c_2,n_0>0, \forall n\ge n_0, 0\le c_1 g(n)\le
f(n)\le c_2 g(n).\]
\[\lim_{n \to \infty} \frac{f(n)}{g(n)}=c\in (0,\infty)\implies f\in \Theta(g).\]

$f$ grows \textcolor{Bittersweet}{as fast as $g$} asymptotically.

\textcolor{Blue}{$\mathbf{f \in o(g)}$} if
\[\forall c>0\quad \exists n_0>0\quad \forall n\ge n_0\quad 0\le f(n)< c\cdot g(n).\]
\[\lim_{n \to \infty} \frac{f(n)}{g(n)}=0\implies f\in o(g).\] 

$f$ grows \textcolor{Bittersweet}{slower than $g$} asymptotically.

\textcolor{Blue}{$\mathbf{f \in \omega(g)}$} if
\[\forall c>0\quad \exists n_0>0\quad \forall n\ge n_0\quad 0\le c\cdot g(n)< f(n).\]
\[\lim_{n \to \infty} \frac{f(n)}{g(n)}=\infty\implies f\in \omega(g).\]

$f$ grows \textcolor{Bittersweet}{faster than $g$} asymptotically.

\paragraph{Properties}
For functions $f,g,h$:
\begin{itemize}
  \item \textcolor{Blue}{Reflexitivty}: $f\in O(f)$, $f\in \Omega(f)$, $f\in \Theta(f)$.
  \item \textcolor{Blue}{Transitivity}: For $A\in \left\{
    O,\Omega,\Theta,o,\omega \right\} $
    \[f\in A(g)\text{ and }g\in A(h)\implies f\in A(h).\]
  \item \textcolor{Blue}{Symmetry}: $f\in \Theta(g)\iff g\in \Theta(f)$.
  \item \textcolor{Blue}{Complementarity}: 
    \[f\in O(g)\iff g\in \Omega(f)\]
    \[f\in o(g)\iff g\in \omega(f)\]
\end{itemize}
